% Abhakseite „Das kann ich“ – Zone nur im Gesamt (\mitzone)
%% TODO Ich-kann-Sätze: Platzhalter wie in den Aufgabendateien
\begin{abhakseite}
\ifdefined\mitzone
\abhakgruppe{Kennst du schon}
\abhak{1}{Brüche kürzen; Bruch ↔ Dezimalzahl ↔ Prozent (Achtelbruch als Dezimalzahl und Prozentsatz)}
\abhak{2}{Anteil einer Menge nehmen (30 \% von 10 Feldern; zwei Drittel von 9 Kugeln)}
\abhak{3}{Relative Häufigkeit berechnen (Anzahl geteilt durch Gesamtzahl)}
\abhak{4}{Anzahl der Zahlen in einem Bereich (101 bis 900 sind 800 Lose; Endziffer 6 in 106 bis 896 sind 80)}
\abhak{5}{Brüche multiplizieren (auch drei Faktoren) und gleichnamige Brüche addieren}
\abhak{6}{Anzahl der Zahlen in einem Bereich (101 bis 900 sind 800 Lose; Endziffer 6 in 106 bis 896 sind 80) – Fehler finden}
\abhak{7}{Anzahl der Zahlen in einem Bereich (101 bis 900 sind 800 Lose; Endziffer 6 in 106 bis 896 sind 80) – selbst rechnen}
\fi
\abhakgruppe{Einheit 1 · Zählen und Ergebnismengen}
\abhak{8}{Zählen}
\abhak{9}{Ergebnisse zu „mindestens einmal …“ aufzählen}
\abhak{10}{Zählen – Fehler finden, Begründen, Darstellungswechsel, Anwendung}
\abhakgruppe{Einheit 2 · Wahrscheinlichkeit einstufig}
\abhak{11}{Wie viele Stufen?}
\abhak{12}{Gleich groß?}
\abhak{13}{Einstufig}
\abhak{14}{Topf mit P = 50 \% auswählen}
\abhak{15}{erwartete Anzahl bei n Versuchen}
\abhak{16}{Einstufig – Fehler finden, Begründen, Darstellungswechsel, Anwendung}
\abhakgruppe{Einheit 3 · Baumdiagramm und Pfadregeln}
\abhak{17}{Ein Pfad oder mehrere?}
\abhak{18}{Nicht oder mindestens?}
\abhak{19}{Mit Zurücklegen}
\abhak{20}{Mit Zurücklegen – Fehler finden, Begründen, Darstellungswechsel, Anwendung}
\abhakgruppe{Einheit 4 · Ohne Zurücklegen}
\abhak{21}{Ohne Zurücklegen}
\abhak{22}{Ohne Zurücklegen – Fehler finden, Begründen, Darstellungswechsel, Anwendung}
\end{abhakseite}
